\documentclass[11pt]{article}
\usepackage[margin=1in]{geometry}
\usepackage{amsmath,amssymb,amsthm,mathtools}
\usepackage{graphicx,microtype,enumitem,xcolor}
\usepackage[colorlinks=true,linkcolor=blue!55!black,citecolor=blue!55!black,urlcolor=blue!55!black]{hyperref}

\newtheorem{theorem}{Theorem}
\newtheorem{remark}{Remark}
\newcommand{\M}{\mathcal M}
\newcommand{\Lcal}{\mathcal L}
\newcommand{\G}{\Gamma}

\title{Stable Exponent Gain for Schur Multipliers\\Toward Coine's Endpoint Problem}
\author{Automated mathematics run \texttt{fa\_banach\_001}, lane 7}
\date{17 August 2026}

\begin{document}
\maketitle

\begin{center}
\fbox{\parbox{0.91\textwidth}{\textbf{Status.} Substantial partial result,
likely valid, pending human review.  The result proves a new exponent-gain
inclusion and gives a direct proof of the known $p=2$ endpoint.  It does not
settle the question for $1<p<2$.}}
\end{center}

\section{Source question and result}

Coine's Problem 5.9 asks whether
\[
                    \M(p,1)=\M(p,p),\qquad 1<p\le 2,
\]
for Schur multipliers on the indicated sequence-space operator classes
\cite{Coine2018}.  The exact question appears on PDF page 24:

\begin{figure}[ht]
\centering
\includegraphics[width=\textwidth]{figures/open_problem_crop.png}
\caption{Problem 5.9 in \cite{Coine2018}.}
\end{figure}

The equality at $p=2$ follows from Bennett's inclusion theorem as quoted by
Coine.  For $1<p<2$, Bennett's equality region reaches every
$\M(p,q)$ with $q<p$ but not the diagonal endpoint $q=p$.  The following
result places a substantial family of diagonal multipliers on the desired
$\M(p,1)$ side.

For an operator $T:E\to F$, write $\gamma_r(T)$ for its $L^r$-factorization
norm.  Write $L_{1,p'}(T)$ for the infimum of $\|S\|\|R\|$ over
factorizations
\[
 E\xrightarrow{R}L^p(\Omega)
  \xrightarrow{I_{p,1}}L^1(\Omega)
  \xrightarrow{S}F
\]
on probability spaces.  Coine's Theorem 4.2 and Corollary 4.5 identify these
operator norms with the relevant Schur multiplier norms.

\begin{theorem}[Stable exponent gain]\label{thm:gain}
Let $1<p<r\le2$.  Every operator
$T:\ell_1\to\ell_\infty$ which factors through an $L^r$ space factors
through a probability-space inclusion $L^p\hookrightarrow L^1$.  Quantitatively,
\[
 L_{1,p'}(T)\le
 \frac{\kappa_{r,p}}{\kappa_{r,1}}\,\gamma_r(T),
 \]
where $\kappa_{r,s}=(\mathbb E|\theta_r|^s)^{1/s}$ and $\theta_r$ is a
normalized symmetric $r$-stable scalar.  Consequently,
\[
                         \boxed{\M(r,r)\subseteq\M(p,1)}.
\]
In particular, $\M(2,2)\subseteq\M(p,1)$ for every $1<p<2$.
At $p=r=2$, with Gaussian constants in place of $\kappa_{r,s}$,
\[
                         \M(2,1)=\M(2,2).
\]
\end{theorem}

\section{Proof intuition}

An $r$-stable random measure turns every vector in an $L^r$ factor into a
random variable whose $L^s$ norm is an exact constant times the original
$L^r$ norm for every $s<r$.  Taking $s=p$ makes the new input map bounded
into $L^p$; taking $s=1$ makes the same map bounded below in $L^1$.  Hence the
old output map is bounded on the resulting subspace when that subspace is
given its $L^1$ norm.  The special range $\ell_\infty$ now matters:
coordinatewise Hahn--Banach extends this output map to all of $L^1$ without
increasing its norm.  This produces exactly the factorization characterizing
$\M(p,1)$.

The strict inequality $p<r$ is structural.  A nondegenerate $r$-stable
variable has finite $s$-moment precisely for $s<r$ (apart from the Gaussian
case $r=2$), so the same mechanism has no $L^p$ endpoint when $r=p<2$.

\section{Proof of Theorem~\ref{thm:gain}}

Fix a factorization
\[
             T=BA,\qquad
 \ell_1\xrightarrow{A}L^r(\Sigma,\nu)
       \xrightarrow{B}\ell_\infty.
\]
Only the closed linear span
\[
 X=\overline{A(\ell_1)}\subseteq L^r(\Sigma,\nu)
\]
is relevant, and $X$ is separable.

For $1<r<2$, the standard symmetric $r$-stable random-measure construction
gives a probability space $(\Omega,\mathbb P)$ and a complex-linear map
$J:X\to L^0(\Omega)$ such that for every $x\in X$, $Jx$ is an isotropic
$r$-stable scalar of scale $\|x\|_r$.  Thus, for every $1\le s<r$,
\begin{equation}\label{eq:moments}
                  \|Jx\|_{L^s}=\kappa_{r,s}\|x\|_r.
\end{equation}
This is the usual stable-process embedding of $L^r$; see
\cite[Chapter 3]{SamorodnitskyTaqqu1994} and the stable-process discussion
used in \cite[Theorem 6.4.19]{AlbiacKalton2006}.  For $r=2$, take the
isonormal complex Gaussian process over the Hilbert space $X$; then
\eqref{eq:moments} holds for every finite $s$, with Gaussian moment
constants.

Since $p<r$, equation \eqref{eq:moments} shows that
\[
 J:X\longrightarrow L^p(\Omega)
 \quad\text{has norm }\kappa_{r,p},
\]
while the same formula with $s=1$ shows that $JX$ is a closed subspace of
$L^1(\Omega)$ and that $J^{-1}:JX\to X$ has norm
$\kappa_{r,1}^{-1}$ when $JX$ carries its inherited $L^1$ norm.

Define
\[
 S_0:JX\longrightarrow\ell_\infty,
 \qquad S_0(Jx)=Bx.
\]
This is well-defined and
\[
             \|S_0\|\le \frac{\|B\|}{\kappa_{r,1}}.
\]
The space $\ell_\infty$ is $1$-injective: writing the coordinate functionals
of $S_0$ and extending each by Hahn--Banach gives an operator
\[
 S:L^1(\Omega)\longrightarrow\ell_\infty
 \quad\text{with}\quad \|S\|=\|S_0\|.
\]
Set $R=JA:\ell_1\to L^p(\Omega)$.  Because the middle map is the canonical
probability-space inclusion $I_{p,1}:L^p\hookrightarrow L^1$, we have
\[
                    T=S I_{p,1}R.
\]
Moreover,
\[
 \|S\|\|R\|
 \le \frac{\kappa_{r,p}}{\kappa_{r,1}}\|B\|\|A\|.
\]
Taking the infimum over all $L^r$ factorizations proves the operator-norm
inequality.

For a matrix symbol $\phi$, Coine's characterization gives
\[
 \|\phi\|_{\M(r,r)}=\gamma_r(u_\phi),
 \qquad
 \|\phi\|_{\M(p,1)}=L_{1,p'}(u_\phi).
\]
The claimed Schur-multiplier inclusion and estimate follow.

When $p=r=2$, use the Gaussian construction, for which both the first and
second moments are finite, to obtain
$\M(2,2)\subseteq\M(2,1)$.  Conversely, a factorization through
$L^2\hookrightarrow L^1$ is in particular a factorization through $L^2$
after composing the last two arrows.  Hence
$\M(2,1)\subseteq\M(2,2)$, completing the endpoint proof. \qed

\section{Verification, limitations, and upgrade audit}

The proof is analytic and has no computational dependency.  Its two
nonformal inputs are standard: the stable random-measure embedding and the
$1$-injectivity of $\ell_\infty$.  All remaining arrows and constants are
displayed explicitly.

This theorem does not prove $\M(p,p)\subseteq\M(p,1)$ for $1<p<2$ and does
not give a counterexample.  The exponent gap cannot simply be closed: a
nondegenerate $p$-stable variable has infinite $p$-th moment.  The associated
universal-symbol reduction also stops short of a contradiction.  If
$\sigma:\ell_1\twoheadrightarrow\ell_p$ is a metric quotient and
$V:\ell_p\hookrightarrow\ell_\infty$ is an isometry, then
$V\sigma\in\Gamma_p$; a hypothetical $(p,1)$ factorization yields only a
factorization of $I_{\ell_p}$ through a subspace of a quotient attached to
$L^p\hookrightarrow L^1$.  It does not yield a direct factorization through
that inclusion.

Eight focused upgrade attempts were made.  Besides the successful stable
gain, they tested: the Hilbert endpoint; the universal identity symbol;
Reisner's $(p,1)$ Gordon--Lewis obstruction; finite-dimensional
quotient/isometry amplification; Hadamard blocks; the source's Rademacher
type estimate at its zero-exponent endpoint; and truncation of $p$-stable
variables.  The detailed log records why none of the remaining routes closes
the question.

\section{Literature and novelty boundary}

Bounded searches on 17 August 2026 covered the run indexes, the exact arXiv
id/title/Problem 5.9 wording, Coine's published paper and thesis, Bennett's
inclusion cases, and combinations of ``$\Gamma_r$'', ``$(p,1)$-factorable'',
and stable-process factorization.  No explicit later resolution of Problem
5.9 and no explicit statement of Theorem~\ref{thm:gain} was located.
Novelty confidence is moderate rather than high because the argument is short
and may be implicit in standard operator-ideal folklore.

Human review is recommended, focusing on the simultaneous complex stable
embedding, the coordinatewise extension into $\ell_\infty$, and the exact
translation through Coine's factorization norms.

\sloppy
\begin{thebibliography}{9}
\bibitem{Coine2018}
C. Coine,
\emph{Schur multipliers on $\mathcal B(L^p,L^q)$},
J. Operator Theory 79 (2018), 301--326; arXiv:1703.08128.

\bibitem{AlbiacKalton2006}
F. Albiac and N. J. Kalton,
\emph{Topics in Banach Space Theory},
Graduate Texts in Mathematics 233, Springer, 2006.

\bibitem{SamorodnitskyTaqqu1994}
G. Samorodnitsky and M. S. Taqqu,
\emph{Stable Non-Gaussian Random Processes},
Chapman \& Hall, 1994.

\bibitem{Reisner1995}
S. Reisner,
\emph{Operators factoring through Banach lattices and ideal norms},
Serdica Math. J. 21 (1995), 75--89.
\end{thebibliography}

\end{document}
